% TOP_PROBLEM: {"problemId": "problem.nearby-lagrangian-conjecture", "canonicalTitle": "Nearby Lagrangian Conjecture", "releaseRank": 270, "primaryDomain": "geometry_topology", "primaryDomainLabel": "Geometry & topology", "displayStatus": "open", "releaseStatus": "open"}
% !TeX program = lualatex
% Definition and sources only; this file is not an LLM solution attempt.
\documentclass[11pt]{article}
\usepackage[margin=25mm]{geometry}
\usepackage{fontspec}
\setmainfont{Latin Modern Roman}
\usepackage{amsmath,amssymb,mathtools}
\usepackage{unicode-math}
\setmathfont{Latin Modern Math}
\usepackage{xurl,hyperref}
\hypersetup{colorlinks=true,urlcolor=blue}
\setcounter{secnumdepth}{0}
\setlength{\parindent}{0pt}
\setlength{\parskip}{5pt}
\setlength{\emergencystretch}{3em}
\begin{document}
\section*{\texorpdfstring{Nearby Lagrangian Conjecture}{Nearby Lagrangian Conjecture}}\label{270}
\subsection{Definitions and mathematical statement}
For the smooth base manifold $N$ in the standard closed-base cotangent formulation, $T^*N$ has canonical one-form $\lambda$ and symplectic form ${\,\mathrm{d}}\lambda$. A closed embedded submanifold $L$ is Lagrangian when $\dim L=\dim N$ and $({\,\mathrm{d}}\lambda)|_L=0$, and exact when $\lambda|_L={\,\mathrm{d}} f$ for some function $f:L\to{\mathbb{R}}$. The conjecture asserts that there is a Hamiltonian isotopy $\Phi_t$ of $T^*N$ with
\[
 \Phi_1(L)=\text{the zero section }N.
\]
The isotopy is generated by a time-dependent Hamiltonian in the source's support convention. Homotopy equivalence of the projection $L\to N$, or abstract diffeomorphism of $L$ and $N$, is weaker than this ambient Hamiltonian-isotopy conclusion.
\par\smallskip\textit{Formulation and concept references:} \href{https://www.math.stonybrook.edu/~fzheng/seminar-f24/notes/week-1-2.pdf}{[S1]}.

\subsection{Short English statement}
Is every closed exact Lagrangian in a cotangent bundle obtainable from its zero section by a Hamiltonian isotopy?
\subsection{Sources}
\begin{itemize}
\item[S1]Formal statement, Conjecture 1.5.
\url{https://www.math.stonybrook.edu/~fzheng/seminar-f24/notes/week-1-2.pdf}.
\textit{Formulation source.}
\end{itemize}
\end{document}
